\begin{table}[!htbp]
\centering
\caption{Main model, HS3, Portuguese Language: ATT by relative period}
\label{tab:ap_event_base_HS3_lp}
\vspace{-0.15cm}
\scriptsize
\setlength{\tabcolsep}{4pt}
\renewcommand{\arraystretch}{0.92}
\resizebox{\textwidth}{!}{%
\begin{tabular}{lcccc}
\toprule
Relative period & TWFE & CS21 & DR-DiD & S\&X approx. \\
\midrule
-4 & 0.062 (0.037) & 0.043 (0.046) & 0.016 (0.046) & 0.044 (0.089) \\
-3 & 0.067 (0.038) & 0.042 (0.040) & 0.014 (0.040) & 0.112 (0.063) \\
-2 & 0.025 (0.032) & 0.024 (0.033) & 0.016 (0.038) & 0.043 (0.045) \\
-1 & 0.000 & 0.000 & 0.000 & 0.000 \\
0 & 0.354 (0.042) & 0.347 (0.043) & 0.358 (0.045) & 0.434 (0.071) \\
1 & 0.451 (0.041) & 0.450 (0.044) & 0.506 (0.046) & 0.541 (0.070) \\
2 & 0.541 (0.042) & 0.541 (0.047) & 0.583 (0.047) & 0.645 (0.076) \\
3 & 0.592 (0.043) & 0.628 (0.049) & 0.682 (0.048) & 0.665 (0.087) \\
4 & 0.595 (0.049) & 0.663 (0.062) & 0.724 (0.058) & 0.722 (0.090) \\
\bottomrule
\end{tabular}
}
\begin{flushleft}
\footnotesize Note: each cell reports the ATT and standard error in parentheses, in standard deviations of annual proficiency. Period \(k=-1\) is normalized as the baseline and therefore appears without a standard error.
\end{flushleft}
\end{table}



